\documentclass[12pt,a4paper]{article}
\usepackage[utf8]{inputenc}
\usepackage[indonesian]{babel}
\usepackage{amsmath,amssymb}
\usepackage{graphicx}
\usepackage{booktabs}
\usepackage{hyperref}
\usepackage{geometry}
\usepackage{float}
\usepackage{caption}
\usepackage{setspace}
\usepackage{natbib}
\geometry{margin=2.5cm}
\onehalfspacing

\title{\textbf{Analisis Pola Makroekonomi dan Aktivitas Keuangan Digital untuk Penguatan Ketahanan Ekonomi Publik: Pendekatan Regresi Huber \textit{First-Differences} pada Prediksi Tingkat Wanprestasi 90 Hari (TWP90) di Indonesia}}

\author{
Tim Riset Datathon METC\\
\textit{Microsoft Datathon Competition 2026}
}
\date{April 2026}

\begin{document}
\maketitle

% ============================================================
% ABSTRACT
% ============================================================
\begin{abstract}
Penelitian ini menganalisis pola makroekonomi dan aktivitas keuangan digital guna merumuskan strategi penguatan ketahanan ekonomi publik melalui prediksi Tingkat Wanprestasi 90 Hari (TWP90) pada sektor \textit{peer-to-peer lending} di 31 provinsi Indonesia. Data panel tahunan periode 2022--2025 digunakan dengan delapan indikator makroekonomi sebagai variabel prediktor, mencakup suku bunga BI, inflasi, PDRB, tingkat pengangguran, penetrasi internet, \textit{Loan-to-Deposit Ratio} (LDR), \textit{Non-Performing Loan} (NPL), dan rasio UMKM. Arsitektur model \textit{First-Differences} (FD) dengan estimator Huber ($\epsilon = 1{,}35$) dipilih untuk mengeliminasi heterogenitas tak terobservasi antar-provinsi dan menangani pencilan (\textit{outlier}) pada data ekonomi regional. Model menghasilkan \textit{Test} $R^2 = 0{,}600$ dan $\text{RMSE} = 0{,}00554$, menunjukkan perbaikan 17{,}9\% dibandingkan estimator OLS standar. Temuan utama mengidentifikasi NPL ($\beta = -0{,}207$) dan inflasi ($\beta = -0{,}201$) sebagai pendorong dominan perubahan TWP90, mencerminkan mekanisme \textit{self-correcting} sektor perbankan. Sistem \textit{decision support} mengklasifikasikan 4 provinsi berisiko tinggi (DKI Jakarta, NTB, Jawa Timur, Jawa Barat) dan 8 provinsi berisiko sedang untuk tahun 2026. Implikasi kebijakan mencakup penguatan pengawasan kredit berbasis risiko regional dan optimalisasi distribusi sumber daya keuangan negara berdasarkan profil risiko provinsi.
\end{abstract}

\textbf{Kata Kunci:} TWP90, regresi Huber, \textit{first-differences}, data panel, sistem peringatan dini, ketahanan ekonomi, \textit{fintech}

% ============================================================
% 1. PENDAHULUAN
% ============================================================
\section{Pendahuluan}

Pertumbuhan pesat sektor keuangan digital di Indonesia, khususnya \textit{peer-to-peer} (P2P) \textit{lending}, telah memberikan akses pembiayaan baru bagi segmen masyarakat yang sebelumnya tidak terjangkau oleh perbankan konvensional. Namun, ekspansi yang cepat ini juga membawa risiko sistemik berupa peningkatan potensi gagal bayar yang dapat mengancam stabilitas keuangan nasional. Tingkat Wanprestasi 90 Hari (TWP90)---didefinisikan sebagai rasio pinjaman yang mengalami keterlambatan pembayaran lebih dari 90 hari terhadap total pinjaman---menjadi indikator kunci kesehatan portofolio kredit digital.

Permasalahan inti yang diangkat dalam penelitian ini adalah: \textit{bagaimana menganalisis pola makroekonomi dan aktivitas keuangan digital untuk merumuskan strategi penguatan ketahanan ekonomi publik dan mengoptimalkan distribusi sumber daya keuangan negara yang tepat sasaran}. Pemahaman terhadap faktor-faktor makroekonomi yang memengaruhi TWP90 pada tingkat provinsi sangat penting bagi pengambil kebijakan untuk merancang intervensi yang bersifat preventif dan berbasis bukti empiris.

Secara spesifik, penelitian ini bertujuan untuk: (1) membangun model ekonometrika panel yang mampu memprediksi perubahan TWP90 antar-provinsi berdasarkan indikator makroekonomi; (2) mengidentifikasi faktor pendorong utama (\textit{key drivers}) risiko kredit digital pada tingkat regional; dan (3) mengembangkan sistem \textit{decision support} yang mengklasifikasikan tingkat risiko provinsi untuk tahun 2026 guna mendukung kebijakan alokasi sumber daya keuangan negara.

% ============================================================
% 2. TINJAUAN PUSTAKA
% ============================================================
\section{Tinjauan Pustaka}

Teori risiko kredit menyatakan bahwa kualitas portofolio pinjaman dipengaruhi oleh kondisi makroekonomi agregat. Hubungan antara variabel-variabel yang digunakan dalam penelitian ini dengan risiko kredit dapat dirangkum sebagai berikut:

\textbf{\textit{Non-Performing Loan} (NPL)} merupakan indikator langsung kualitas aset perbankan. Secara teori, peningkatan NPL menandakan deteriorasi kualitas kredit. Namun, dalam kerangka dinamis, kenaikan NPL dapat memicu respons pengetatan standar kredit oleh lembaga keuangan, yang secara paradoks memperbaiki kualitas debitur baru.

\textbf{Inflasi} memengaruhi kapasitas pembayaran debitur melalui erosi daya beli riil. Inflasi yang tinggi dapat meningkatkan tekanan finansial pada rumah tangga, namun juga dapat mendorong kebijakan moneter kontraktif yang membatasi ekspansi kredit.

\textbf{\textit{Loan-to-Deposit Ratio} (LDR)} mengukur agresivitas penyaluran kredit relatif terhadap dana pihak ketiga. LDR yang tinggi mengindikasikan ekspansi kredit yang agresif, yang berpotensi menurunkan standar penyaringan debitur.

\textbf{Tingkat Pengangguran Terbuka (TPT)} berkorelasi dengan kemampuan pembayaran debitur; peningkatan pengangguran umumnya berkorelasi positif dengan gagal bayar.

\textbf{Penetrasi Internet} merupakan proksi inklusi keuangan digital. Peningkatan akses internet memperluas basis pengguna \textit{fintech}, termasuk segmen dengan literasi keuangan rendah yang berpotensi meningkatkan risiko gagal bayar.

\textbf{PDRB} mencerminkan kapasitas ekonomi regional. Pertumbuhan ekonomi yang ekspansif dapat meningkatkan aktivitas pinjaman, sementara kontraksi ekonomi meningkatkan tekanan gagal bayar.

% ============================================================
% 3. METODOLOGI
% ============================================================
\section{Metodologi}

\subsection{Data dan Variabel}

Penelitian ini menggunakan data panel dari 31 provinsi di Indonesia untuk periode 2022--2025 ($N = 31$, $T = 4$). Data bersumber dari basis data perbankan nasional yang memuat variabel bulanan yang kemudian diagregasi menjadi rata-rata tahunan. Total observasi pada level tahunan adalah 124 baris.

Variabel dependen adalah $Y_{it}$ = TWP90 rata-rata tahunan provinsi $i$ pada tahun $t$. Variabel independen ($\mathbf{X}_{it}$) terdiri dari delapan indikator:

\begin{table}[H]
\centering
\caption{Deskripsi Variabel Independen}
\label{tab:variabel}
\begin{tabular}{lll}
\toprule
\textbf{Simbol} & \textbf{Variabel} & \textbf{Satuan} \\
\midrule
$x_1$ & Suku Bunga BI (\textit{BI Rate}) & Persen \\
$x_2$ & Inflasi & Rasio desimal \\
$x_3$ & $\ln(\text{PDRB})$ & Log rupiah \\
$x_4$ & Tingkat Pengangguran Terbuka & Persen \\
$x_5$ & Penetrasi Internet & Rasio (0--1) \\
$x_6$ & \textit{Loan-to-Deposit Ratio} & Rasio \\
$x_7$ & \textit{Non-Performing Loan} & Rasio desimal \\
$x_8$ & Rasio UMKM & Rasio \\
\bottomrule
\end{tabular}
\end{table}

Tahun 2021 dieksklusi dari analisis karena variabel penetrasi internet ($x_5$) tidak tersedia. Analisis awal menunjukkan bahwa imputasi linier untuk $x_5$ pada tahun 2021 menghasilkan degradasi signifikan pada performa model (\textit{Test} $R^2$ turun dari $0{,}40$ menjadi $-2{,}9$), sehingga keputusan eksklusi bersifat konservatif dan metodologis.

\subsection{Pra-pemrosesan Data}

\begin{enumerate}
\item \textbf{Agregasi Temporal:} Data bulanan diagregasi menjadi rata-rata tahunan per provinsi untuk setiap variabel. Hal ini dilakukan karena variabel makroekonomi bersifat \textit{slow-moving} dan tidak menunjukkan variasi bermakna pada frekuensi bulanan.
\item \textbf{Transformasi Logaritmik:} PDRB ditransformasi menggunakan logaritma natural ($\ln$) untuk menormalkan distribusi yang sangat miring (\textit{right-skewed}) dan menginterpretasikan koefisien sebagai elastisitas.
\item \textbf{Penanganan Nilai Hilang:} Variabel $x_5$ pada observasi awal ditangani dengan \textit{backward fill} per provinsi.
\item \textbf{Standardisasi:} Seluruh fitur distandardisasi menggunakan \textit{StandardScaler} ($z = \frac{x - \mu}{\sigma}$) sebelum estimasi untuk memastikan komparabilitas koefisien dan stabilitas numerik.
\end{enumerate}

\subsection{Arsitektur Model: \textit{First-Differences} (FD)}

Model panel konvensional dengan \textit{fixed effects} (LSDV) memiliki formulasi:
\begin{equation}
Y_{it} = \alpha_i + \boldsymbol{\beta}' \mathbf{X}_{it} + \epsilon_{it}
\end{equation}
di mana $\alpha_i$ menangkap efek tetap provinsi. Dengan $N = 31$ provinsi dan hanya $T = 4$ periode waktu, estimasi LSDV membutuhkan 39 parameter (31 \textit{dummy} + 8 koefisien) dari 93 observasi diferensi, menghasilkan rasio observasi-per-parameter yang sangat rendah dan menyebabkan matriks singular.

Pendekatan \textit{First-Differences} mengeliminasi $\alpha_i$ secara matematis dengan mengambil selisih antar-periode:
\begin{equation}
\Delta Y_{it} = Y_{it} - Y_{i,t-1} = \boldsymbol{\beta}' (\mathbf{X}_{it} - \mathbf{X}_{i,t-1}) + (\epsilon_{it} - \epsilon_{i,t-1})
\end{equation}
yang disederhanakan menjadi:
\begin{equation}
\Delta Y_{it} = \boldsymbol{\beta}' \Delta \mathbf{X}_{it} + \Delta \epsilon_{it}
\end{equation}

Transformasi ini menghilangkan seluruh karakteristik provinsi yang tidak berubah terhadap waktu (\textit{time-invariant unobserved heterogeneity}) tanpa memerlukan estimasi parameter tambahan.

\subsection{Estimator Huber (\textit{M-Estimator})}

Analisis eksploratori menunjukkan bahwa distribusi residual memiliki ekor tebal (\textit{heavy-tailed}), terutama disebabkan oleh provinsi-provinsi dengan volatilitas TWP90 ekstrem (misalnya NTB dan DKI Jakarta). Estimator OLS standar, yang meminimalkan \textit{sum of squared residuals}, sangat sensitif terhadap pencilan tersebut.

Estimator Huber merupakan \textit{M-estimator} yang meminimalkan fungsi kerugian campuran:
\begin{equation}
\min_{\boldsymbol{\beta}} \sum_{i,t} \rho_\epsilon \left( \frac{\Delta Y_{it} - \boldsymbol{\beta}' \Delta \mathbf{X}_{it}}{\sigma} \right)
\end{equation}
dengan fungsi kerugian Huber:
\begin{equation}
\rho_\epsilon(r) = \begin{cases}
\frac{1}{2} r^2 & \text{jika } |r| \leq \epsilon \\
\epsilon |r| - \frac{1}{2} \epsilon^2 & \text{jika } |r| > \epsilon
\end{cases}
\end{equation}

Parameter $\epsilon = 1{,}35$ dipilih berdasarkan konvensi statistik yang memberikan efisiensi 95\% relatif terhadap OLS pada distribusi normal, namun secara substansial lebih robust terhadap pencilan.

\subsection{Strategi Evaluasi}

Data diferensi dibagi secara temporal:
\begin{itemize}
\item \textbf{Training:} $\Delta$2023 dan $\Delta$2024 (62 observasi)
\item \textbf{Testing:} $\Delta$2025 (31 observasi)
\end{itemize}

Prediksi level-space direkonstruksi: $\hat{Y}_{i,2025} = Y_{i,2024} + \widehat{\Delta Y}_{i,2025}$.

Metrik evaluasi yang digunakan:
\begin{align}
\text{RMSE} &= \sqrt{\frac{1}{n}\sum_{i=1}^{n}(Y_i - \hat{Y}_i)^2} \\
\text{MAE} &= \frac{1}{n}\sum_{i=1}^{n}|Y_i - \hat{Y}_i| \\
R^2 &= 1 - \frac{\sum(Y_i - \hat{Y}_i)^2}{\sum(Y_i - \bar{Y})^2}
\end{align}

Sebagai \textit{benchmark}, digunakan model naif \textit{random walk}: $\hat{Y}_{i,t} = Y_{i,t-1}$.

% ============================================================
% 4. HASIL DAN PEMBAHASAN
% ============================================================
\section{Hasil dan Pembahasan}

\subsection{Performa Model}

Tabel~\ref{tab:metrics} menyajikan perbandingan performa model pada \textit{level-space}.

\begin{table}[H]
\centering
\caption{Perbandingan Metrik Evaluasi Model}
\label{tab:metrics}
\begin{tabular}{lcccc}
\toprule
\textbf{Model} & \textbf{$n$} & \textbf{RMSE} & \textbf{MAE} & \textbf{$R^2$} \\
\midrule
Huber FD (Train) & 62 & 0,00588 & 0,00359 & 0,641 \\
\textbf{Huber FD (Test)} & \textbf{31} & \textbf{0,00554} & \textbf{0,00431} & \textbf{0,600} \\
Naif $y_t = y_{t-1}$ & 31 & 0,00477 & 0,00320 & 0,704 \\
\midrule
OLS FD (Test, sebelumnya) & 31 & 0,00675 & --- & 0,406 \\
\bottomrule
\end{tabular}
\end{table}

Model Huber FD menunjukkan perbaikan substansial dibandingkan OLS standar: RMSE turun 17,9\% (dari 0,00675 menjadi 0,00554) dan $R^2$ meningkat dari 0,406 menjadi 0,600. Gap \textit{overfitting} antara train dan test hanya 6,1\%, mengindikasikan stabilitas generalisasi yang baik. Tidak terdapat prediksi negatif (0/31 provinsi), menunjukkan validitas fisik model.

Perlu dicatat bahwa model masih berada di bawah \textit{benchmark} naif dalam hal presisi titik (RMSE 0,00554 vs 0,00477). Hal ini merupakan konsekuensi dari keterbatasan jumlah observasi training ($n = 62$). Namun, model berhasil mengungguli naif pada 14 dari 31 provinsi dan memberikan nilai tambah substantif dalam hal interpretabilitas koefisien dan kemampuan simulasi skenario.

\subsection{Koefisien dan \textit{Driver} Ekonomi}

Tabel~\ref{tab:coef} menyajikan koefisien model pada skala asli, diurutkan berdasarkan nilai absolut.

\begin{table}[H]
\centering
\caption{Koefisien Regresi Huber FD (Skala Asli)}
\label{tab:coef}
\begin{tabular}{lcl}
\toprule
\textbf{Variabel ($\Delta$)} & \textbf{Koefisien} & \textbf{Interpretasi} \\
\midrule
NPL ($x_7$) & $-0{,}207$ & Kenaikan NPL $\rightarrow$ penurunan TWP90 \\
Inflasi ($x_2$) & $-0{,}201$ & Kenaikan inflasi $\rightarrow$ penurunan TWP90 \\
Penetrasi Internet ($x_5$) & $+0{,}017$ & Digitalisasi $\rightarrow$ kenaikan TWP90 \\
LDR ($x_6$) & $+0{,}013$ & Ekspansi kredit $\rightarrow$ kenaikan TWP90 \\
$\ln$(PDRB) ($x_3$) & $+0{,}011$ & Pertumbuhan ekonomi $\rightarrow$ kenaikan TWP90 \\
BI Rate ($x_1$) & $-0{,}003$ & Suku bunga naik $\rightarrow$ TWP90 turun \\
UMKM ($x_8$) & $-0{,}003$ & Dampak marginal kecil \\
TPT ($x_4$) & $-0{,}001$ & Dampak marginal kecil \\
\bottomrule
\end{tabular}
\end{table}

\textbf{Temuan 1: NPL sebagai \textit{driver} dominan dengan arah negatif.} Koefisien negatif NPL ($\beta = -0{,}207$) menunjukkan bahwa kenaikan 1 poin persentase pada $\Delta$NPL berasosiasi dengan penurunan 0,207 poin persentase pada $\Delta$TWP90. Temuan ini bersifat kontra-intuitif namun konsisten dengan teori \textit{credit rationing}: ketika NPL meningkat, lembaga keuangan secara proaktif mengetatkan standar penyaluran kredit baru, menghasilkan \textit{pool} debitur yang lebih berkualitas sehingga tingkat gagal bayar pada periode berikutnya menurun. Mekanisme ini merupakan bukti empiris dari \textit{self-correcting mechanism} pada sektor keuangan.

\textbf{Temuan 2: Inflasi berasosiasi negatif.} Koefisien inflasi ($\beta = -0{,}201$) mengindikasikan bahwa kenaikan inflasi berasosiasi dengan penurunan TWP90. Dalam konteks Indonesia pasca-pandemi, kenaikan inflasi seringkali diikuti oleh pengetatan moneter (kenaikan BI Rate) yang membatasi ekspansi kredit baru, sehingga secara tidak langsung menurunkan proporsi pinjaman berisiko tinggi.

\textbf{Temuan 3: Digitalisasi meningkatkan risiko marginal.} Koefisien positif penetrasi internet ($\beta = +0{,}017$) dan LDR ($\beta = +0{,}013$) menunjukkan bahwa perluasan akses keuangan digital dan ekspansi kredit berkorelasi dengan peningkatan TWP90, meskipun dengan magnitude yang lebih kecil dibandingkan NPL dan inflasi.

\subsection{Klasifikasi Risiko Provinsi 2026}

Sistem \textit{decision support} (A3) melatih ulang model Huber FD pada seluruh data (93 observasi diferensi) dan menghasilkan klasifikasi risiko berdasarkan TWP90 aktual 2025 dan prediksi skenario \textit{stress}:

\begin{table}[H]
\centering
\caption{Klasifikasi Risiko Provinsi untuk Tahun 2026}
\label{tab:risk}
\begin{tabular}{lcl}
\toprule
\textbf{Level} & \textbf{Jumlah} & \textbf{Provinsi} \\
\midrule
HIGH & 4 & DKI Jakarta (4,67\%), NTB (4,04\%), \\
 & & Jawa Timur (3,41\%), Jawa Barat (3,36\%) \\
MEDIUM & 8 & DIY, Sumsel, Jateng, Lampung, \\
 & & Sumbar, Banten, Kalsel, Jambi \\
LOW & 19 & Sisanya \\
\bottomrule
\end{tabular}
\end{table}

Kriteria klasifikasi: HIGH jika TWP90 aktual $> 3\%$ atau TWP90 \textit{stress} $> 4\%$; MEDIUM jika TWP90 aktual $> 2\%$ atau TWP90 \textit{stress} $> 3\%$; LOW untuk sisanya.

\subsection{Analisis Skenario}

Tiga skenario disimulasikan untuk tahun 2026:
\begin{itemize}
\item \textbf{Baseline:} $\Delta X = 0$ untuk seluruh variabel (TWP90 stabil).
\item \textbf{Stress:} NPL $+1$ ppt, inflasi $+0{,}5\%$, TPT $+0{,}3$, LDR $-2\%$.
\item \textbf{Optimistik:} NPL $-0{,}5$ ppt, pertumbuhan PDRB $+5\%$.
\end{itemize}

Pada skenario \textit{stress}, model memprediksi penurunan TWP90 sebesar 4,8\%--28,6\% secara relatif, yang mencerminkan mekanisme pengetatan kredit oleh lembaga keuangan sebagai respons terhadap tekanan ekonomi. Temuan ini konsisten dengan koefisien negatif NPL yang telah dibahas.

% ============================================================
% 5. IMPLIKASI KEBIJAKAN
% ============================================================
\section{Implikasi Kebijakan}

Berdasarkan temuan empiris, tiga strategi kebijakan dirumuskan untuk menjawab permasalahan penguatan ketahanan ekonomi publik dan optimalisasi distribusi sumber daya keuangan negara:

\textbf{Strategi 1: Pengawasan Kredit Berbasis Risiko Regional.} Klasifikasi risiko provinsi menunjukkan konsentrasi risiko TWP90 tinggi pada empat provinsi utama (DKI Jakarta, NTB, Jawa Timur, Jawa Barat). Otoritas pengawas keuangan (OJK) disarankan mengimplementasikan kerangka pengawasan berjenjang: intensifikasi audit dan pembatasan ekspansi pinjaman baru pada provinsi berisiko tinggi, serta peningkatan cadangan CKPN (\textit{Cadangan Kerugian Penurunan Nilai}) untuk platform \textit{fintech} yang beroperasi di wilayah tersebut.

\textbf{Strategi 2: Optimalisasi Alokasi Dana Penjaminan Kredit.} Model menunjukkan bahwa variabel NPL dan inflasi merupakan \textit{leading indicator} yang efektif untuk memprediksi arah perubahan TWP90. Pemerintah dapat mengoptimalkan distribusi dana penjaminan kredit (misalnya melalui Jamkrindo atau Askrindo) dengan mengalokasikan proporsi dana yang lebih besar kepada provinsi-provinsi yang terklasifikasi HIGH dan MEDIUM, guna memitigasi dampak sistemik dari potensi gagal bayar. Pendekatan berbasis data ini memastikan sumber daya keuangan negara terdistribusi secara tepat sasaran.

\textbf{Strategi 3: Penguatan Literasi Keuangan di Wilayah Ekspansi Digital.} Koefisien positif penetrasi internet terhadap TWP90 mengindikasikan bahwa perluasan akses keuangan digital tanpa disertai peningkatan literasi keuangan berpotensi meningkatkan risiko gagal bayar. Program edukasi keuangan digital perlu diprioritaskan di provinsi-provinsi dengan pertumbuhan penetrasi internet yang tinggi namun memiliki literasi keuangan yang rendah, sebagai strategi preventif penguatan ketahanan ekonomi publik.

% ============================================================
% 6. KESIMPULAN
% ============================================================
\section{Kesimpulan}

Penelitian ini berhasil membangun sistem peringatan dini (\textit{early warning system}) untuk TWP90 pada sektor P2P \textit{lending} di Indonesia menggunakan arsitektur \textit{First-Differences} dengan estimator Huber. Model menjelaskan 60\% varians TWP90 pada data uji dengan perbaikan RMSE sebesar 17,9\% dibandingkan estimator OLS standar, dan gap \textit{overfitting} yang rendah (6,1\%).

Temuan kunci menunjukkan bahwa NPL dan inflasi merupakan pendorong utama perubahan TWP90 dengan arah negatif---mencerminkan mekanisme pengetatan kredit yang bersifat \textit{self-correcting}. Sistem \textit{decision support} yang dihasilkan mengklasifikasikan 31 provinsi ke dalam tiga tingkat risiko dan menyediakan simulasi skenario untuk mendukung pengambilan kebijakan berbasis bukti.

\subsection{Limitasi}

\begin{enumerate}
\item \textbf{Ukuran sampel terbatas:} Agregasi tahunan menghasilkan hanya 62 observasi training, membatasi kemampuan model menangkap pola kompleks.
\item \textbf{Asosiasi, bukan kausalitas:} Koefisien mencerminkan korelasi parsial, bukan hubungan kausal. Inferensi kebijakan yang lebih kuat memerlukan pendekatan variabel instrumental.
\item \textbf{Regime spesifik:} Data mencakup periode pasca-pandemi (2022--2025). Perubahan struktural ekonomi di masa depan dapat mengubah hubungan antar-variabel.
\item \textbf{Presisi titik:} Model masih di bawah \textit{benchmark} naif dalam hal RMSE absolut, sehingga lebih tepat digunakan untuk perankingan dan analisis arah daripada prediksi numerik presisi.
\end{enumerate}

\subsection{Rekomendasi Penelitian Lanjutan}

Peningkatan granularitas data ke frekuensi kuartalan akan secara substansial meningkatkan \textit{degrees of freedom} dan daya prediktif model. Eksplorasi interaksi antar-variabel (misalnya $\Delta\text{NPL} \times \Delta\text{LDR}$) dan penggunaan estimator panel dinamis (Arellano-Bond GMM) dapat memperdalam pemahaman mekanisme transmisi risiko kredit.

\end{document}
